%\begin{name}
%	{Biên soạn: Thầy Nguyễn Thành Nhân  \& Phản biện: Thầy Nguyễn Cường}
%	{Đề chuẩn hóa}
%\end{name}

\caulc
\Opensolutionfile{ans}[ans/ans-cau-TN-1]

\begin{ex}%[De-chuan-hoa-so-1]%[Nguyễn Thành Nhân]%[1D6H4-2]
	Tập nghiệm của phương trình $\log_2x=\log_2(2x+1)$ là 
	\choice
	{$\{-1\}$}
	{\True $\varnothing$}
	{$\{0\}$}
	{$\{1\}$}
	\loigiai{
		Ta có 
		$\log_2x=\log_2(2x+1)\Leftrightarrow \heva{&x>0\\&x=2x+1}\Leftrightarrow\heva{&x>0\\&x=-1}\Rightarrow \text{phương trình vô nghiệm}.$
	}
\end{ex}

\begin{ex}%[De-chuan-hoa-so-1]%[Nguyễn Thành Nhân]%[1D7N2-1]
Mệnh đề nào sau đây {\bf đúng}?
\choice%26
{$\left(\sin^2{2x}\right)^\prime = \sin4x$}
{$\left(\cos{2x}\right)^\prime = 2\sin2x$}
{\True $\left(\cos^2{2x}\right)^\prime = -2\sin4x$}
{$\left(\sin{2x}\right)^\prime =-2 \cos2x$}

\loigiai{
Ta có
$$\left(\sin^2{2x}\right)^\prime=2\sin 2x\cdot (\sin2x)^\prime=2\sin 2x\cdot 2\cos 2x=2\sin 4x.$$
$$\left(\cos{2x}\right)^\prime=-\sin 2x\cdot (2x)^\prime=-2\sin 2x.$$
$$\left(\cos^2{2x}\right)^\prime = 2\cos2x\cdot (\cos2x)^\prime =-2\cdot 2\cos2x \cdot \sin2x=-2\sin4x.$$
$$\left(\sin{2x}\right)^\prime=\cos 2x\cdot (2x)^\prime=2\cos 2x.$$
}
\end{ex}
\begin{ex}%[De-chuan-hoa-so-1]%[Nguyễn Thành Nhân]%[1D7H2-1]
Hàm số $y=(x+1)\sqrt{x^2+1}$ có đạo hàm là
\choice%30
{\True $\dfrac{2x^2+x+1}{\sqrt{x^2+1}}$}
{$\dfrac{x^2+x+1}{\sqrt{x^2+1}}$}
{$\dfrac{x^2+2x+1}{\sqrt{x^2+1}}$}
{$\dfrac{x^2+x+2}{\sqrt{x^2+1}}$}
\loigiai{
$y^\prime =\sqrt{x^2+1}+(x+1)\cdot \dfrac{x}{\sqrt{x^2+1}}=\dfrac{x^2+1+x^2+x}{\sqrt{x^2+1}}=\dfrac{2x^2+x+1}{\sqrt{x^2+1}}$.}
\end{ex}
\begin{ex}%[De-chuan-hoa-so-1]%[Nguyễn Thành Nhân]%[1D7H3-2]
	Cho hàm số $y = x \sin x$. Hệ thức nào sau đây đúng?
	\choice
	{$y'' - y = 2 \sin x$}
	{$y'' + y = 2 \sin x$}
	{$y'' - y = 2 \cos x$}
	{\True $y'' + y = 2 \cos x$}
	\loigiai{
	Ta có: $y' = \sin x + x \cos x$. Và
	$y''= \cos x + \cos x - x \sin x = 2 \cos x - x \sin x$. \\
	Vậy $y'' + y = 2 \cos x - x \sin x + x \sin x = 2 \cos x$.
	}
\end{ex}
\begin{ex}%[De-chuan-hoa-so-1]%[Nguyễn Thành Nhân]%[1D2H3-5]
	Cho hàm số $f(x) = \dfrac{1}{5}x^5 + \dfrac{4}{3}x^3 - 5x + 3$. Mệnh đề nào sau đây \textbf{sai}?
	\choice
	{Phương trình $f(x) = 0$ có nghiệm trên khoảng $(-1;1)$}
	{Phương trình $f(x) = 0$ có nghiệm trên khoảng $(0;+ \infty)$}
	{\True Hàm số đã cho gián đoạn tại $x_0 = \dfrac{1}{5}$}
	{Hàm số $f(x)$ liên tục trên $\mathbb{R}$}
	\loigiai{
	Tập xác định: $\mathscr{D} = \mathbb{R}$.
		Ta có $f\left(\dfrac{1}{2} \right) \cdot f\left(\dfrac{4}{5} \right) < 0 \Rightarrow \exists x_0 \in \left(0;\dfrac{9}{10}\right): f(x_0) =0$. \\
	Do đó: \lq\lq Phương trình $f(x) = 0$ có nghiệm trên khoảng $(-1;1)$\rq\rq\, và \lq\lq Phương trình $f(x) = 0$ có nghiệm trên khoảng $(0;+ \infty)$ \rq\rq\, là hai mệnh đề đúng.\\
	Hàm số $f(x)$ là hàm đa thức nên liên tục trên tập xác định $\mathscr{D} = \mathbb{R}$.\\
	Nên mệnh đề \lq\lq Hàm số đã cho gián đoạn tại $x_0 = \dfrac{1}{5}$\rq\rq \, là mệnh đề sai.
	}
\end{ex}
\begin{ex}%[De-chuan-hoa-so-1]%[Nguyễn Thành Nhân]%[1H8H1-3]
Cho hình chóp $S.ABCD$ có đáy là hình vuông cạnh $a$, cạnh bên $SA$ vuông góc với mặt đáy, $SA=a$, gọi $M$ là trung điểm $SB$. Góc giữa $AM$ và $BD$ bằng
\choice%19
{$90^\circ$}
{$30^\circ$}
{\True $60^\circ$}
{$45^\circ$}
\loigiai{\immini{Gọi $N$ lần lượt là trung điểm của $SD$. Khi đó $MN$ là dường trung bình của tam giác $SBD$. Ta có $$\Rightarrow MN\parallel BD \Rightarrow (AM,BD)=(AM,MN).$$
Ta có $AM$ và $AN$ lần lượt là đường trung tuyến của các tam giác vuông $SAB$ và $SAD$ nên $AM=\dfrac{SB}{2}=\dfrac{SD}{2}=AN=\dfrac{a\sqrt{2}}{2}$. Lại có $MN$ là đường trung bình tam giác $SBD$ nên $MN=\dfrac{BD}{2}=\dfrac{a\sqrt{2}}{2}$.\\
  Từ đó suy ra tam giác $AMN$ đều cạnh $\dfrac{a}{\sqrt{2}}$. Vậy $(AM,MN)=60^\circ$.}
{\begin{tikzpicture}[>=stealth,line join=round,line cap=round,font=\footnotesize,scale=0.5]
\clip (-4,-4) rectangle (9,7);
	\def\r{6.5};
	\def\g{120};
	\def\gC{30};
	\path
	(0,0) coordinate (A)
	(-3,-3) coordinate (B)
	(5,-3) coordinate (C)
	(8,0) coordinate (D)
	(0,6) coordinate (S)
	;
	\draw  (S)--(B)--(C)--(D)--(S)--(C);
	\coordinate (M) at ($(S)!0.5!(B)$);
	\coordinate (N) at ($(S)!0.5!(D)$);
	%\path[name path=d1] (A)--(C) ;% Gán tên đường $AC$ là $d1$.
	%\path[name intersections={of=d1 and d2}] (intersection-1) coordinate (I);%Tìm giao của hai đường $d1$, $d2$ đặt tên điểm là I$.
	\draw[dashed]  (S)--(A)--(D)(B)--(A)--(M)circle (1.5pt) node[above left]{$M$}--(N)circle (1.5pt) node[above right]{$N$}--(A) (B)--(D);
	\fill[black] (A) circle (1.5pt) node[below right]{$A$}(C) circle (1.5pt) node[below right]{$C$}(B) circle (1.5pt) node[below left]{$B$}(D) circle (1.5pt) node[above right]{$D$}(S) circle (1.5pt) node[above right]{$S$};
	\draw pic[draw,angle radius = 5pt] {right angle =S--A--B};%Vẽ góc vuông BKC
		\draw pic[draw,angle radius = 5pt] {right angle =S--A--D};
\end{tikzpicture}
}
}
\end{ex}

\begin{ex}%[De-chuan-hoa-so-1]%[Nguyễn Thành Nhân]%[1H8N4-1]
Mệnh đề nào sau đây đúng?
\choice%22
{Đường thẳng $d$ là đường vuông góc chung của hai đường thẳng nhau $a$, $b$ khi và chỉ khi $d$ vuông góc với cả $a$ và $b$}
{\True Cho hai mặt phẳng vuông góc với nhau, nếu một đường thẳng nằm trong mặt phẳng này và vuông góc với giao tuyến của hai mặt phẳng thì vuông góc với mặt phẳng kia}
{Nếu đường thẳng $a$ và mặt phẳng $(P)$ cùng vuông góc với đường thẳng $d$ thì đường thẳng $a$ song song với mặt phẳng $(P)$}
{Nếu hai đường thẳng phân biệt cùng vuông góc với một đường thẳng thứ ba thì chúng song song với nhau}
\loigiai{
Hai mặt phẳng vuông góc với nhau, nếu một đường thẳng nằm trong mặt phẳng này và vuông góc với giao tuyến của hai mặt phẳng thì vuông góc với mặt phẳng kia.}
\end{ex}

\begin{ex}%[De-chuan-hoa-so-1]%[Nguyễn Thành Nhân]%[1D7N2-3]
Cho parabol $(P):y=x^2-3x$. Tiếp tuyến của $(P)$ đi qua điểm $A(5;10)$ có phương trình là
\choice%25
{$y=5x-15$}
{\True $y=7x-25$}
{$y=x+5$}
{$y=3x-5$}
\loigiai{
Ta thấy điểm $A(5;10)$ thuộc parabol nên $A$ là tiếp điểm.\\
$y^\prime = 2x-3\Rightarrow y^\prime (5)=7$.\\
Phương trình tiếp tuyến là $y=7x-25$.}
\end{ex}

\begin{ex}%[De-chuan-hoa-so-1]%[Nguyễn Thành Nhân]%[1H8N6-4]
Cho hình chóp tam giác đều $S.ABC$ có cạnh đáy bằng $2a$, góc giữa mặt bên và mặt đáy bằng $60^\circ$. Độ dài cạnh bên của hình chóp bằng
\choice%27
{\True $\dfrac{a\sqrt{21}}{3}$}
{$\dfrac{a\sqrt{7}}{3}$}
{$\dfrac{a\sqrt{17}}{3}$}
{$\dfrac{a\sqrt{15}}{3}$}
\loigiai{
\immini{Gọi $H$ là chân đường cao, suy ra $H$ là tâm tam giác đều $ABC$.\\
Gọi $M$ là trung điểm $BC$ suy ra góc giữa mặt bên $(SBC)$ và $(ABC)$ là góc $\widehat{SMH}=60^\circ$.\\
Ta có $AH=\dfrac{2}{3}AM=\dfrac{2a\sqrt{3}}{3}; HM=\dfrac{1}{3}AM=\dfrac{a\sqrt{3}}{3}$.\\
$SH=HM\cdot \tan \widehat{SMH}=\dfrac{a\sqrt{3}}{3}\cdot \sqrt{3}=a$.\\
Xét $\triangle SHA: SA^2=SH^2+AH^2=\dfrac{4a^2}{3}+a^2=\dfrac{7a^2}{3}\Rightarrow SA=\dfrac{a\sqrt{21}}{3}$.}
{\begin{tikzpicture}[>=stealth,line join=round,line cap=round,font=\footnotesize,scale=0.5]
\clip (-1,-5) rectangle (12,7);
	\def\r{6.5};
	\def\g{120};
	\def\gC{30};
	\path
	(0,0) coordinate (A)
	(7,-4) coordinate (B)
	(11,0) coordinate (C)
	(6,6) coordinate (S)
	;
	\draw  (S)--(A)--(B)--(C)--(S)--(B);
	\coordinate (M) at ($(B)!0.5!(C)$);
	\coordinate (H) at ($(A)!0.666!(M)$);
	%\path[name path=d1] (A)--(C) ;% Gán tên đường $AC$ là $d1$.
	%\path[name intersections={of=d1 and d2}] (intersection-1) coordinate (I);%Tìm giao của hai đường $d1$, $d2$ đặt tên điểm là I$.
	\draw[dashed]  (S)--(H)--(M)(C)--(A)--(H)--(B)(H)--(C)(M)--(S);
	\fill[black] (A) circle (1.5pt) node[below left]{$A$}(C) circle (1.5pt) node[below right]{$C$}(B) circle (1.5pt) node[below left]{$B$}(S) circle (1.5pt) node[above right]{$S$}(H) circle (1.5pt) node[below left]{$H$}(M) circle (1.5pt) node[below right]{$M$};
	\draw pic[draw,angle radius = 5pt] {right angle =S--H--B};%Vẽ góc vuông BKC
		\draw pic[draw,angle radius = 5pt] {right angle =S--H--C};
\end{tikzpicture}
}}
\end{ex}

\begin{ex}%[De-chuan-hoa-so-1]%[Nguyễn Thành Nhân]%[1H8H5-3]
	Cho tứ diện $ABCD$ có các tam giác $ABC$, $DBC$ vuông cân và nằm trong hai mặt phẳng vuông góc với nhau, $AB = AC = DB = DC = 2a$. Tính khoảng cách từ $B$ đến mặt phẳng $(ACD)$.
	\choice
	{\True $\dfrac{2a\sqrt{6}}{3}$}
	{$\dfrac{2a\sqrt{3}}{3}$}
	{$a\sqrt{6}$}
	{$\dfrac{a\sqrt{6}}{2}$}
	
	\loigiai{
	\immini{
	Gọi $E$ là trung điểm cạnh $BC$. Tam giác $ABC$ vuông cân tại $A$ nên $AE$ vừa là trung tuyến vừa là đường cao. \\
	Suy ra: $AE \perp BC \Rightarrow AE \perp (BCD)$. \\
	Gọi $K$ là trung điểm của cạnh $DC$. Ta có\\ $\heva{EK \parallel BD \\ BD \perp DC} \Rightarrow DC \perp EK$ \hfill(1)\\
	Lại có $AE \perp DC$ \hfill(2)\\
	Từ (1) và (2), suy ra $DC \perp (AEK) \Rightarrow (AEK) \perp (ACD)$ theo giao tuyến $AK$.\\
	}
		{
	\begin{tikzpicture}[>=stealth,line join=round,line cap=round,font=\footnotesize,scale=0.5]
\clip (-1,-5) rectangle (13,7);
	\def\r{6.5};
	\def\g{120};
	\def\gC{30};
	\path
	(0,0) coordinate (B)
	(6,-4) coordinate (C)
	(12,0) coordinate (D)
	(3,6) coordinate (A)
	;
	\coordinate (E) at ($(B)!0.5!(C)$);
	\coordinate (H) at ($(A)!0.666!(M)$);
	\coordinate (K) at ($(C)!0.5!(D)$);
	\coordinate (H) at ($(A)!(E)!(K)$);
	\draw  (A)--(B)--(C)--(D)--(A)--(C)(A)--(E)(A)--(K);
	%\path[name path=d1] (A)--(C) ;% Gán tên đường $AC$ là $d1$.
	%\path[name intersections={of=d1 and d2}] (intersection-1) coordinate (I);%Tìm giao của hai đường $d1$, $d2$ đặt tên điểm là I$.
	\draw[dashed]  (B)--(D)--(E)--(K) circle (1.5pt) node[below]{$K$}(E)--(H)circle (1.5pt) node[above right]{$H$};
	\fill[black] (A) circle (1.5pt) node[above left]{$A$}(C) circle (1.5pt) node[below right]{$C$}(B) circle (1.5pt) node[below left]{$B$}(S) circle (1.5pt) node[above right]{$S$}(D) circle (1.5pt) node[below right]{$D$}(E) circle (1.5pt) node[below left]{$E$};
	\draw pic[draw,angle radius = 5pt] {right angle =A--E--D};%Vẽ góc vuông BKC
	\draw pic[draw,angle radius = 5pt] {right angle =A--H--E};
\end{tikzpicture}
	}
	\noindent
	Trong mặt phẳng $(AEK)$ kẻ $EH \perp AK$ với $H \in AK$. Suy ra $EH \perp (ACD)$.\\
	Do đó $\mathrm{d}(B, (ACD)) = 2 \mathrm{d}(E, (ACD)) = 2EH$. \\
	Ta có $AE \cdot BC = AB \cdot AC \Rightarrow AE = \dfrac{AB \cdot AC}{BC} = a \sqrt{2}$. $EK = \dfrac{1}{2}BD = a$. \\
	Tam giác $AEK$ vuông tại $E$ có $EH$ là đường cao. \\
	$EH^2 = \dfrac{AE^2 \cdot EK^2}{AE^2 + EK^2} = \dfrac{2a^2 \cdot a^2}{2a^2 + a^2} = \dfrac{2a^2}{3} \Leftrightarrow EH = \dfrac{a \sqrt{6}}{3}$. Nên $\mathrm{d}(B, (ACD)) = \dfrac{2a \sqrt{6}}{3}$.
	}
\end{ex}
\begin{ex}%[1-Đề HKI, SGDDT - KonTum, 2020-2021]%[Trần Phú Hiếu]%[1D2G5-4]
	Hai xạ thủ độc lập bắn vào một mục tiêu. Xác suất trúng mục tiêu của xạ thủ thứ nhất là $0,7$. Xác suất trúng mục tiêu của xạ thủ thứ hai là $0,8$. Xác suất để mục tiêu bị bắn trúng là
	\choice
	{\True $P=0{,}94$}
	 {$P=0{,}56$}
	{$P=0{,}08$}
	{ $P=0{,}06$}
	\loigiai{Gọi $A_1$ là biến cố: \lq\lq người thứ nhất bắn trúng\rq\rq. \\
		Gọi $A_2$ là biến cố: \lq\lq người thứ hai bắn trúng\rq\rq.\\
		Suy ra $\mathrm{P}(A)=0{,}8$;  $\mathrm{P}=0{,}7$.\\
		Gọi $B$ là biến cố: \lq\lq mục tiêu bị bắn trúng\rq\rq.\\
		Ta có 
		\[\mathrm{P}(B)=\mathrm{P}(A_1)\cdot \mathrm{P}(A_2)+\mathrm{P}(A_1)\cdot \mathrm{P}\left(\overline{A_2}\right)+\mathrm{P}(A_2)\cdot \mathrm{P}\left(\overline{A_1}\right)=0{,}8\cdot 0{,}7+0{,}8\cdot 0{,}3+0{,}2\cdot 0{,}7=0{,}94.\]
}
\end{ex}
\begin{ex}%[De-chuan-hoa-so-1]%[Nguyễn Thành Nhân]%[1D7H2-1]
	Cho hàm số $f(x) = \dfrac{x^3}{3} - \dfrac{x^2}{2} - 12x + \dfrac{1}{3}$. Tập nghiệm $\mathscr{T}$ của bất phương trình $f'(x) \leq 0$ là
	\choice
	{$\mathscr{T} \left(- \infty; - 3 \right] \cup \left[4; + \infty \right)$}
	{$\mathscr{T} = \varnothing$}
	{\True $\mathscr{T} =\left[-3;4 \right]$}
	{$\mathscr{T} = \left(- \infty; + \infty \right)$}
	\loigiai{
	Tập xác định: $\mathscr{D} = \mathbb{R}$. \\
	Ta có: $f'(x) = x^2 - x - 12$. \\
	$f'(x) \leq 0 \Leftrightarrow -3 \leq x \leq 4$.
	}
\end{ex}
\Closesolutionfile{ans}
\indapan{7}{ans/ans-cau-TN-1}

\cauds
%\setcounter{ex}{0}
\Opensolutionfile{ans}[ans/ans-cau-DS-1]
\begin{ex}%[De-chuan-hoa-so-1]%[Nguyễn Thành Nhân]%[1D7V2-2]
	Cho hàm số $y=f(x)=\sqrt{x+\sqrt{x^2+1}}$. Khi đó
	\choiceTF
	{\True Hàm số xác định với mọi $x\in \mathbb{R}$}
	{\True $\lim\limits _{x\to -\sqrt{3}}f(x)=\sqrt{2-\sqrt{3}}$}
	{$f'(0)=\dfrac{1}{\sqrt{2}}$}
	{\True $2\sqrt{x^2+1} \cdot y'-y=0$}
	\loigiai{
		\begin{itemchoice}
			\itemch Đúng.\\
			 Vì $\sqrt{x^2+1}>\left|x\right|\geq -x\Rightarrow \sqrt{x^2+1}+x>0,\,\forall x\in \mathbb{R}$ nên hàm số có tập xác định là $\mathbb{R}$.
			\itemch Đúng.\\
			Ta có $\lim\limits _{x\to -\sqrt{3}}f(x)=\sqrt{-\sqrt{3}+\sqrt{(-\sqrt{3})^2+1}}=\sqrt{2-\sqrt{3}}$.
			\itemch Sai.\\
			Vì $f'(x)=\dfrac{1+\tfrac{x}{\sqrt{x^2+1}}}{2\sqrt{x+\sqrt{x^2+1}}}\Rightarrow f'(0)=\dfrac{1}{2}$.
			\itemch Đúng.\\
			Vì $y'=\dfrac{1+\tfrac{x}{\sqrt{x^2+1}}}{2\sqrt{x+\sqrt{x^2+1}}}=\dfrac{\sqrt{x^2+1}+x}{\sqrt{x^2+1}\cdot 2\sqrt{x+\sqrt{x^2+1}}}=\dfrac{y^2}{2y\sqrt{x^2+1}}=\dfrac{y}{2\sqrt{x^2+1}}$.
	 Nên $2\sqrt{x^2+1} \cdot y'-y=0$. 
		\end{itemchoice}
	}	
\end{ex}
\begin{ex}%%[De-chuan-hoa-so-1]%[Nguyễn Thành Nhân]%[1D6V2-2]
	Cho hàm số $y = x + \dfrac{4}{x - 3}$ có đồ thị $(C)$ và đường thẳng $d\colon y = -3x + 2$. Khi đó
	\choiceTF
	{Điểm $N(-1;0)$ thuộc đồ thị $(C)$}
	{Đồ thị $(C)$ và đường thẳng $d$ cắt nhau tại hai điểm có tích hoành độ giao điểm bằng $-\dfrac{5}{2}$}
	{\True Tiếp tuyến của đồ thị $(C)$ tại điểm có hoành độ bằng $1$ song song với trục hoành}
	{\True $M_1(2;-2),M_2(4;8)$ là hai điểm thuộc đồ thị $(C)$ mà tiếp tuyến của $(C)$ tại đó song song với đường thẳng $d$}
	\loigiai{
	\begin{itemchoice}
	\itemch Sai.\\
	Với $x=-1$ thì $y=-2$ nên điểm $N(-1;0)$ không thuộc đồ thị $(C)$.
	\itemch Sai\\
	Phương trình hoành độ giao điểm của $(C)$ và $(d)$ là 
	\begin{eqnarray*}
	&& x + \dfrac{4}{x - 3}=-3x+2\\
	&\Leftrightarrow & \heva{&x\ne 3\\&4x-2+\dfrac{4}{x-3}=0}\\
	&\Leftrightarrow & \heva{& x\ne 3\\& (4x-2)(x-3)+4=0}\\
	&\Leftrightarrow & \heva{&x\ne 3\\&2x^2-7x+5=0.}
		\end{eqnarray*}
	Phương trình này có $\Delta=9>0$ nên có hai nghiệm. Theo định lý Viète tích hai nghiệm bằng $\dfrac{5}{2}$.
	\itemch Đúng.\\
	Với $x=1$ thì $y=-1$ nên điểm $K(1;-1)$ không thuộc trục $Ox$.\\
	Lại có $y'=1-\dfrac{4}{(x-3)^2}$ nên $y'(1)=0$. Do đó tiếp tuyến của $(C)$ tại $K(1;-1)$ song song với trục hoành.
	\itemch Đúng.\\
	Đặt $f(x)=x + \dfrac{4}{x - 3}$.
	Ta có $f'(x) = 1 - \dfrac{4}{\left(x - 3 \right)^2}$. \\
	Gọi $M\left(x_0; f\left(x_0 \right) \right)$ là tọa độ tiếp điểm. \\
	Hệ số góc của tiếp tuyến tại điểm $\left(x_0; f\left(x_0 \right) \right)$ là $y'(x_0) = \dfrac{x^2_0 - 6x_0 + 5}{\left(x_0 - 3 \right)^2}$. \\
	Vì tiếp tuyến song song với đường thẳng $d:y = -3x + 2 $ nên 
	$$y'(x_0) = - 3 \Leftrightarrow \heva{&4x^2_0 - 24x_0 + 32 = 0 \\ &x_0 \neq 3} \Leftrightarrow \hoac{&x_0 = 4 \\ &x_0 = 2.}$$
	Với $x_0=2$, suy ra tiếp điểm $M_1(2;-2)$ và $f'(2)=-3$. Phương trình tiếp tuyến tại $M_1$ là 
	\[y+2=-3(x-2)\Leftrightarrow y=-3x+4.\]
	Với $x_0=4$, suy ra tiếp điểm $M_2(4;8)$ và $f'(4)=-3$. Phương trình tiếp tuyến tại $M_1$ là 
	\[y-8=-3(x-4)\Leftrightarrow y=-3x+20.\]
	Cả hai tiếp tuyến đều song song với đường thẳng $d$. Suy ra có hai tiếp điểm thỏa yêu cầu bài $M_1(2;-2),M_2(4;8)$.
	\end{itemchoice}
	}
\end{ex}

\begin{ex}%[De-chuan-hoa-so-1]%[Nguyễn Thành Nhân]%[1H8H7-3]
	Cho hình chóp tứ giác đều $S.ABCD$ có tất cả các cạnh bằng nhau và bằng $a$. Khi đó
	\choiceTF
	{\True $(SAC)\perp (SBD)$}
	{\True Các cạnh bên tạo với đáy một góc $45^{\circ}$}
	{\True Côsin của góc giữa mặt bên và mặt đáy bằng $\dfrac{\sqrt{3}}{3}$}
	{Thể tích khối chóp $S.ABCD$ bằng $\dfrac{a^3\sqrt{2}}{2}$ }
	\loigiai{
	\immini{
	\begin{itemchoice}
	\itemch Đúng\\
	 Gọi $O$ là tâm của $ABCD$. Khi đó $SO\perp (ABCD)$ nên $SO\perp BD$. \\
	 Lại có $BD\perp AC$ nên $BD\perp (SAC)$. Suy ra $(SBD)\perp (SAC)$.
	\itemch Đúng\\
	Góc giữa cạnh bên và đáy bằng góc $\widehat{SCO}$.\\
	Ta có $SC=a$, $OC=\dfrac{a\sqrt{2}}{2}$. Tam giác $SOC$ vuông ở $O$ nên 
	\[\cos \widehat{SCO}=\dfrac{OC}{SC}=\dfrac{\sqrt{2}}{2}\Rightarrow \widehat{SCO}=45^{\circ}.\]
	\itemch Đúng\\
	Gọi $M$ là trung điểm của $CD$. Ta có 
	$$\heva{&OM\perp CD\\&SO\perp CD}\Rightarrow CD\perp (SOM)\Rightarrow CD\perp SM.$$
	Do đó góc giữa mặt bên và mặt đáy bằng góc $\widehat{SMO}$.\\
	Tam giác $SCD$ đều cạnh $a$ nên $SM=\dfrac{a\sqrt{3}}{2}$. Do đó
	\[\cos \widehat{SMO}=\dfrac{OM}{SM}=\dfrac{\frac{a}{2}}{\frac{a\sqrt{3}}{2}}=\dfrac{\sqrt{3}}{3}.\]
	\itemch Sai\\
	Ta có $SO=OC=\dfrac{a\sqrt{2}}{2}$. Thể tích khối chóp $S.ABCD$ là
	\[V=\dfrac{1}{3}\cdot SO\cdot S_{ABCD}=\dfrac{1}{3}\cdot \dfrac{a\sqrt{2}}{2}\cdot a^2=\dfrac{a^3\sqrt{2}}{6}.\]
	\end{itemchoice}}
	{\begin{tikzpicture}[>=stealth,line join=round,line cap=round,font=\footnotesize,scale=0.4]
\clip (-5,-5) rectangle (11,7);
	\def\r{6.5};
	\def\g{120};
	\def\gC{30};
	\path
	(0,0) coordinate (A)
	(-4,-4) coordinate (B)
	(6,-4) coordinate (C)
	(10,0) coordinate (D)
	(3,6) coordinate (S)
	(3,-2) coordinate (O)
	;
	\coordinate (M) at ($(D)!0.5!(C)$);
	\draw  (S)--(B)--(C)--(D)--(S)--(C)(M)--(S);
	%\path[name path=d1] (A)--(C) ;% Gán tên đường $AC$ là $d1$.
	%\path[name intersections={of=d1 and d2}] (intersection-1) coordinate (I);%Tìm giao của hai đường $d1$, $d2$ đặt tên điểm là I$.
	\draw[dashed]  (B)--(A)--(S)(B)--(D)--(A)--(C)(S)--(O)circle (1.5pt) node[below]{$O$}--(M)circle (1.5pt) node[below]{$M$};
	\fill[black] (A) circle (1.5pt) node[above left]{$A$}(C) circle (1.5pt) node[below right]{$C$}(B) circle (1.5pt) node[below left]{$B$}(S) circle (1.5pt) node[above right]{$S$}(D) circle (1.5pt) node[below right]{$D$};
	%\draw pic[draw,angle radius = 5pt] {right angle =A--E--D};
\end{tikzpicture}}
	}
\end{ex}

\begin{ex}%[De-chuan-hoa-so-1]%[Nguyễn Thành Nhân]%[1D7H1-1]
	Cho hàm số $f(x) = \heva{&(m-1)x^2 + m &\text{khi} \ x \geq 2 \\&3x -1 &\text{khi}\ x < 2}(m \text{ là tham số)}$. Khi đó
	\choiceTF
	{\True Hàm số xác định với mọi $x\in \mathbb{R}$}
	{$f(2)=5$}
	{Với $m = 2$ thì hàm số liên tục trên $\mathbb{R}$}
	{\True Hàm số không có đạo hàm tại $x=2$}
	\loigiai{
	\begin{itemchoice}
	\itemch Đúng\\
	Với $x\geq 2$ thì $f(x)=(m-1)x^2 + m $ là hàm số đa thức nên xác định. Với $x<2$ thì $f(x)=3x-1$ là hàm số đa thức nên xác định. Vậy hàm số xác định với mọi $x\in \mathbb{R}$.
	\itemch Sai\\
	 Ta có $f(2)=4(m-1)+m=5m-4$.
	\itemch Sai\\
	Để hàm số liên tục trên $\mathbb{R}$ thì chỉ cần hàm số liên tục thêm tại $x=2$. Do đó
	\[\lim \limits_{x\to 2^+}f(x)=\lim \limits_{x\to 2^-}f(x)=f(2)\Leftrightarrow 4(m-1)+m=5\Leftrightarrow m=\dfrac{9}{5}.\] 
	\itemch Đúng\\
	Để hàm số có đạo hàm tại $x=2$ thì hàm số phải liên tục tại $x=2$. Do đó $m=\dfrac{9}{5}$.\\
	 Với $m=\dfrac{9}{5}$ thì $f(x)=\heva{&\tfrac{4}{5}x^2+\tfrac{9}{5}&\text{khi} \ x \geq 2\\&3x -1 &\text{khi}\ x < 2}$. \\
	 Ta có $f'(2^+)=\dfrac{16}{5}$ và $f'(2^-)=3$. Vì $f'(2^+)\ne f'(2^-)$ nên hàm số không có đạo hàm tại $x=2$.
	\end{itemchoice}
	}
\end{ex}
\Closesolutionfile{ans}
\indapan{4}{ans/ans-cau-DS-1}

\caukq
\Opensolutionfile{ans}[ans/ans-cau-KQ-1]
\begin{ex}%[De-chuan-hoa-so-1]%[Nguyễn Thành Nhân]%[1D7H3-1]
Cho hàm số $f(x)=\sin^2x-x^2+1$. Tính $f^{\prime \prime}\left(\dfrac{\pi}{2}\right)$.
\shortans[]{$-4$}
\loigiai{
$f^\prime (x)=2\sin x \cdot \cos x-2x=\sin2x -2x$.\\
$f^{\prime \prime}(x)=2\cos2x-2$. Vậy 
$f^{\prime \prime}\left(\dfrac{\pi}{2}\right)=2\cos \pi -2=-4$.
}
\end{ex}
\begin{ex}%[De-chuan-hoa-so-1]%[Nguyễn Thành Nhân]%[1D7H3-1]
	Có bao nhiêu cặp số nguyên dương $(a; b)$ thỏa mãn $\log_{a}b+6 \log_{b} a=5$ và $a \in[2 ; 2005]$, $b \in[2;2005]$?
	\shortans[]{$54$}
	\loigiai{ 
		Với các điều kiện trên, ta có
		\begin{eqnarray*}
			&&\log_{a}b+6 \log_{b} a=5
			\\
			&\Leftrightarrow& \log_a b +6 \dfrac{1}{\log_a b}=5
			\\
			&\Leftrightarrow& \log^2_a b -5\log_a b+6=0
			\\
			&\Leftrightarrow& \hoac{& \log_a b =2 \\ & \log_a b =3}
			\\
			&\Leftrightarrow& \hoac{&  b=a^2 \\ & b=a^3.}
		\end{eqnarray*}
		Khi $b=a^2 \Rightarrow 2\le a^2 \le 2005 \Rightarrow \sqrt 2 \le a \le \sqrt{2005}$.
		\\
		Do $a$ nguyên nên $2\le a \le 44$, vậy có $44-2+1=43$ số $a$ thoả mãn nên có $43$ cặp số $a$, $b$ thoả mãn yêu cầu bài toán.
		\\
		Khi $b=a^3 \Rightarrow 2\le a^3 \le 2005 \Rightarrow \sqrt[3] 2 \le a \le \sqrt[3]{2005}$.
		\\
		Do $a$ nguyên nên $2\le a \le 12$, vậy có $12-2+1=11$ số $a$ thoả mãn nên có $11$ cặp số $a$, $b$ thoả mãn yêu cầu bài toán.
		\\
		Vậy có tổng cộng $43+11=54$ cặp số thoả mãn.
	} 
\end{ex} 
\begin{ex}%%[De-chuan-hoa-so-1]%[Nguyễn Thành Nhân]%[1d7h2-1]
	Cho hàm số $f(x) = \dfrac{2x + m - 1}{x + m}\ (m \text{ là tham số})$. Tìm  giá trị nguyên không dương của tham số $m$ để $f'(x) >0,\ \forall x \neq -m$.
	\shortans[]{$0$}
	\loigiai{
	Tập xác định của hàm số là  $\mathscr{D} = \mathbb{R} \setminus \{-m \}$. \\
	Ta có $f'(x) = \dfrac{m + 1}{\left(x + m \right)^2}$. Do đó
	$$f'(x) >0,\ \forall x \neq -m \Leftrightarrow \heva{&m > -1 \\ & x \neq - m} \Leftrightarrow m >-1.$$
	Vậy $m=0$.
	}
\end{ex}
\begin{ex}%%[De-chuan-hoa-so-1]%[Nguyễn Thành Nhân]%[1D7V2-7]
	Cho biểu thức $P(x)=(x-1)^{2024}=a_0+a_1x+a_2x^2+\cdots+a_{2024}x^{2024}$. Tính giá trị của biểu thức $$S=a_1+2a_2+3a_3+\cdots+2024a_{2024}.$$
	\shortans[]{$0$}
	\loigiai{
	Đạo hàm hai vế của biểu thức $P(x)$, ta được
	\[2024(x-1)^{2023}=a_1+2a_2x+\cdots+2024a_{2024}x^{2023}.\quad (1)\]
	Trong $(1)$, thay $x=1$ vào hai vế ta được
	\[0=a_1+2a_2+3a_3+\cdots+2024a_{2024}.\]
	Vậy $S=0$.
	}
\end{ex}
\begin{ex}%[De-chuan-hoa-so-1]%[Nguyễn Thành Nhân]%[1H8N6-4]
Cho hình chóp tam giác đều $S.ABC$ có cạnh đáy bằng $2\sqrt{21}$, góc giữa mặt bên và mặt đáy bằng $60^\circ$. Tính độ dài cạnh bên của hình chóp.
\shortans[]{$7$}
\loigiai{
\immini{Đặt $a=\sqrt{21}$. Gọi $H$ là chân đường cao, suy ra $H$ là tâm tam giác đều $ABC$.\\
Gọi $M$ là trung điểm $BC$ suy ra góc giữa mặt bên $(SBC)$ và $(ABC)$ là góc $\widehat{SMH}=60^\circ$.\\
Ta có $AH=\dfrac{2}{3}AM=\dfrac{2a\sqrt{3}}{3}; HM=\dfrac{1}{3}AM=\dfrac{a\sqrt{3}}{3}$.\\
$SH=HM\cdot \tan \widehat{SMH}=\dfrac{a\sqrt{3}}{3}\cdot \sqrt{3}=a$.\\
Xét $\triangle SHA: SA^2=SH^2+AH^2=\dfrac{4a^2}{3}+a^2=\dfrac{7a^2}{3}\Rightarrow SA=\dfrac{a\sqrt{21}}{3}=7$.}{\begin{tikzpicture}[>=stealth,line join=round,line cap=round,font=\footnotesize,scale=0.5]
\clip (-1,-5) rectangle (12,8);
	\def\r{6.5};
	\def\g{120};
	\def\gC{30};
	\path
	(0,0) coordinate (A)
	(4,-4) coordinate (B)
	(10,0) coordinate (C);
	\coordinate (M)at ($(B)!0.5!(C)$);
	\coordinate (H) at ($(A)!0.666!(M)$);
	\coordinate (S') at ($(A)!(H)!(C)$);
	\coordinate (S) at ($(H)!6!(S')$);
	\draw  (S)--(A)--(B)--(C)--(S)--(B)(M)--(S);
	\draw[dashed]  (A)--(H)--(S)(A)--(C)--(H)--(M);
	\fill[black] (A) circle (1.5pt) node[above left]{$A$}(C) circle (1.5pt) node[below right]{$C$}(B) circle (1.5pt) node[below left]{$B$}(S) circle (1.5pt) node[above right]{$S$}(M) circle (1.5pt) node[right]{$M$}(H) circle (1.5pt) node[below]{$H$};
	\draw pic[draw,angle radius = 5pt] {right angle =S--H--A};
	\draw pic[draw,angle radius = 5pt] {right angle =A--M--B};
\end{tikzpicture}
}}
\end{ex}
\begin{ex}%[De-chuan-hoa-so-1]%[Nguyễn Thành Nhân]%[1H8V5-4]
	Cho hình chóp $S.ABCD$ có đáy là hình thoi cạnh $4$, góc $\widehat{ABC}=60^\circ$, cạnh bên $SA=SB=SC$, mặt bên $(SCD)$ tạo với mặt đáy góc $60^\circ$. Tính khoảng cách giữa $AB$ và $SD$.
	\shortans[]{$2$}
	\loigiai
	{
		\immini{Đặt cạnh hình thoi là $a=4$. Tam giác $ABC$ có $AB=AC$ và $\widehat{ABC}=60^{\circ}$ nên là tam giác đều. Vì $SA=SB=SC$ nên hình chiếu $H$ của đỉnh $S$ lên $(ABCD)$ là tâm đường tròn ngoại tiếp tam giác $ABC$, cũng là trọng tâm. \\
		Ta có $CH\perp AB$ nên $CH\perp CD$, $CD\perp SH$. Suy ra $CD\perp (SHC)$. Suy ra $CD\perp SC$ và $CD\perp CH$.\\
		Do đó góc giữa mặt bên $(SCD)$ và mặt đáy là $\widehat{SCH}=60^{\circ}$.\\
		Vì $AB\parallel CD$ nên $AB\parallel (SCD)$.
		}
		{\begin{tikzpicture}[>=stealth,line join=round,line cap=round,font=\footnotesize,scale=0.4]
\clip (-6,-5) rectangle (11,7);
	\def\r{6.5};
	\def\g{120};
	\def\gC{30};
	\path
	(0,0) coordinate (A)
	(-5,-4) coordinate (B)
	(6,-4) coordinate (C)
	(10,0) coordinate (D)
	(3,-2) coordinate (O)
	;
	\coordinate (H) at ($(B)!0.666!(O)$);
	\coordinate (S') at ($(A)!(H)!(D)$);
	\coordinate (S) at ($(H)!3!(S')$);
	\draw  (S)--(B)--(C)--(D)--(S)--(C);
	\coordinate (K) at ($(A)!0.5!(B)$);
	\coordinate (M) at ($(S)!(H)!(C)$);
	\draw[dashed]  (B)--(A)--(S)(B)--(D)--(A)--(C)(S)--(H)circle (1.5pt) node[below]{$H$}--(C)--(K) circle (1.5pt) node[above]{$K$}(H)--(M) circle (1.5pt) node[right]{$M$};
	%\path[name path=d1] (A)--(C) ;% Gán tên đường $AC$ là $d1$.
	%\path[name intersections={of=d1 and d2}] (intersection-1) coordinate (I);%Tìm giao của hai đường $d1$, $d2$ đặt tên điểm là I$.
	
	\fill[black] (A) circle (1.5pt) node[above left]{$A$}(C) circle (1.5pt) node[below right]{$C$}(B) circle (1.5pt) node[below left]{$B$}(S) circle (1.5pt) node[above right]{$S$}(D) circle (1.5pt) node[below right]{$D$}(O) circle (1.5pt) node[below]{$O$};
	\draw pic[draw,angle radius = 5pt] {right angle =S--H--D};
	\draw pic[draw,angle radius = 5pt] {right angle =H--M--C};
\end{tikzpicture}
		}
		Gọi $K$ là trung điểm của $AB$ thì $KC=\dfrac{3}{2}HC$. Kẻ $HM\perp SC\Rightarrow HM\perp (SCD)$, do đó
		\[\mathrm{d}\left(AB,SD\right)=\mathrm{d}\left(AB,(SCD)\right)=\mathrm{d}\left(K,(SCD)\right)=\dfrac{3}{2}\mathrm{d}\left(H,(SCD)\right)=\dfrac{3}{2}HM.\]
		Ta có $CH=\dfrac{a\sqrt{3}}{3}$, $SH=CH\cdot \tan 60^{\circ}=\dfrac{a\sqrt{3}}{3}\cdot \sqrt{3}=a$.
		Tam giác $SHM$ vuông tại $H$ có $HM$ là đường cao nên 
		\[HM=\dfrac{SH\cdot HC}{\sqrt{CH^2+HC^2}}=\dfrac{a\cdot \dfrac{a\sqrt{3}}{3}}{\sqrt{a^2+\dfrac{a^2}{3}}}=\dfrac{a}{2}=2.\]
	}
\end{ex}

\Closesolutionfile{ans}
\indapan{7}{ans/ans-cau-KQ-1}